% ---- en-chunk1.tex ----
\label{VI}
\section{Introduction}

Contrary to what had been announced in
the introduction to the preceding expos\'e, it has
turned out to be impossible to do descent in the category
of preschemes, even in special cases, without having
first developed with sufficient care the language of
descent in general categories.

The notion of ``descent'' provides the general framework for all
procedures of ``gluing'' of objects, and
consequently of ``gluing'' of categories. The most
classical case of gluing is relative to the datum of a
topological space $X$ and a covering of~$X$ by open sets $X_i$; if
one is given for every $i$ a fiber space (say) $E_i$ over
$X_i$, and for every pair $(i,j)$ an isomorphism $f_{ji}$ of
$E_i|X_{ij}$ onto~$E_j|X_{ij}$ (where one sets $X_{ij}=X_i\cap X_j$),
satisfying a well-known transitivity condition (which one
writes in abbreviated form $f_{kj}f_{ji}=f_{ki}$), one
knows that there exists a fiber space $E$ over~$X$, defined up to
isomorphism by the condition that one have isomorphisms
$f_i\colon E|X_i\isomto E_i$, satisfying the relations
$f_{ji}=f_jf_i^{-1}$ (with the usual abuse of writing). Let $X'$
be the sum space of the $X_i$, which is thus a fiber space over
$X$ (\ie equipped with a continuous map $X'\to X$). One can
interpret more concisely the datum of the $E_i$
as a fiber space $E'$ over~$X'$, and the datum of the $f_{ji}$
as an isomorphism between the two inverse images (by the two
canonical projections) $E''_1$ and $E''_2$ of~$E'$ on~$X''=X'\times_X
X'$, the gluing condition then being able to be written as an
identity between isomorphisms of fiber spaces $E'''_1$ and
$E'''_3$ over the triple fibered product $X'''=X'\times_X X'\times_X
X'$ (where $E'''_i$ denotes the inverse image of~$E'$ on~$X'''$
by the canonical projection of index $i$). The construction of~$E$,
starting from~$E'$ and from~$f$, is a typical case
% original p. 146
of a procedure of ``descent''. Moreover, starting from a fiber
space $E$ over~$X$, one says that $X$ is ``locally trivial'', of
fiber $F$, if there exists an open covering $(X_i)$ of~$X$ such that
the $E|X_i$ are isomorphic to $F\times X_i$, or what amounts
to the same, such that the inverse image $E'$ of~$E$ on~$X'=\coprod_i
X_i$ is isomorphic to $X'\times F$.

Thus, the notion of ``gluing'' of objects like that of
``localization'' of a property, are tied to
the study of certain types of ``changes of base'' $X'\to X$.
In Algebraic Geometry, many other types of change
of base, and notably the faithfully flat morphisms $X'\to X$,
must be regarded as corresponding to a
procedure of ``localization'' relative to
preschemes, or other objects, ``over''~$X$. This type of
localization is used just as much as simple
topological localization (which is moreover a special case of it). The same holds
(to a lesser extent) in Analytic Geometry.

Most of the proofs, reducing to
verifications, are omitted or merely sketched; if need
be we specify the less evident diagrams
which arise in a proof.

\section{Universes, categories, equivalence of categories}
\label{VI.1}

To avoid certain logical difficulties, we shall admit
here the notion of \emph{Universe}, which is a set ``large enough''
that one does not leave it by the usual operations of
set theory; an ``\emph{axiom of Universes}'' guarantees
that every object lies in a Universe. For details, see
a book in preparation by C\ptbl Chevalley and the speaker\footnote{%
The definitive authors are C\ptbl Chevalley and P\ptbl Gabriel. The book is to appear in the year 3000. In the meantime, \cf also SGA~4 I.}. Thus, the symbol~$\Ens$
designates, not the category of all sets (a notion
which has no meaning), but the category of sets which
lie in a given Universe (which we shall not specify
here in the notation). Likewise, $\Cat$
will designate the category of categories lying in
the Universe in question, the ``morphisms'' from an object $X$ of~$\Cat$
to another $Y$ being by definition the \emph{functors}
from~$X$ to $Y$.

If
% original p. 147
$\cal{C}$ is a category, we denote by $\Ob(\cal{C})$
\emph{the set of objects} of~$\cal{C}$, by~$\Fl(\cal{C})$
\emph{the set of arrows} of~$\cal{C}$ (or morphisms of
$\cal{C}$). We shall thus write $X\in\Ob(\cal{C})$, avoiding
the common abuse of notation $X\in\cal{C}$. If $\cal{C}$ and $\cal{C}'$
are two categories, a \emph{functor} from~$\cal{C}$ to
$\cal{C}'$ will always be what is commonly called a
\emph{covariant} functor from~$\cal{C}$ to $\cal{C}'$; its datum implies
that of the target category and the source
category, $\cal{C}$ and $\cal{C}'$. The functors from~$\cal{C}$ to
$\cal{C}'$ form a set, denoted $\Hom(\cal{C},\cal{C}')$, which
is the set of objects of a category denoted
$\SheafHom(\cal{C},\cal{C}')$.
By definition, a \emph{contravariant functor} from~$\cal{C}$ to
$\cal{C}'$ is a functor from the \emph{opposite category}
$\cal{C}^\circ$
of~$\cal{C}$ to~$\cal{C}'$.

We shall admit the notion of \emph{projective limit} and of
\emph{inductive limit} of a functor $F\colon \cal{I}\to \cal{C}$, and
in particular the most common special cases of this notion:
cartesian products and fibered products, dual notions of
direct sums and of amalgamated sums, and the usual
formal properties of these operations.

For example, in the category $\Cat$ introduced above, the
projective limits (relative to categories $\cal{I}$
lying in the chosen Universe) exist; the set of objects
(\resp the set of arrows) of the projective limit
category $\cal{C}$ of the $\cal{C}_i$, is obtained by taking the
projective limit of the sets of objects (\resp of the sets of arrows)
of the categories $\cal{C}_i$. The best-known case is that of the
product of a family of categories. We shall constantly use
in the sequel the fibered product of two categories over a
third.

For everything concerning categories and functors, while
awaiting the book in preparation already mentioned, see
\cite{VI.1} (which is necessarily very incomplete, even as regards
the generalities sketched in the
present no.).

Let us take this opportunity to make explicit the notion of equivalence of
categories,
which is not presented in a satisfactory way
in~\cite{VI.1}. A functor $F\colon\cal{C}\to \cal{C}'$ is said to be
\emph{faithful}
(\resp \emph{fully faithful})
if for every pair of objects $S$, $T$ of~$\cal{C}$, the map
$u\mto F(u)$ from~$\Hom(S,T)$ to $\Hom(F(S),F(T))$ is injective
(\resp bijective). One says that $F$ is an \emph{equivalence} of
categories if
% original p. 148
$F$ is fully faithful, and if moreover every object $S'$ of
$\cal{C}'$ is isomorphic to an object of the form $F(S)$. One shows
that this amounts to the same as saying that there exists a functor $G$ from
$\cal{C}'$ to $\cal{C}$ \emph{quasi-inverse of} $F$, \iev such that
$GF$ is isomorphic to $\id_{\cal{C}}$. When this is so, the
datum of a functor $G\colon\cal{C}'\to \cal{C}$ and of an
isomorphism $\varphi\colon GF\to \id_{\cal{C}'}$ is equivalent to
the datum, for every $S'\in\Ob(\cal{C}')$, of a pair $(S,u)$
formed of an object $S$ of~$\cal{C}$ and an isomorphism $u\colon
F(S)\to S'$, namely $(G(S),\varphi(S))$. (With these notations, there exists
a unique functor $\cal{C}'\to\cal{C}$ having the given
map $S\mto G(S)$ as object-map, and such that the map
$S\mto \varphi(S)$ is a homomorphism of functors $FG\to
\id_{\cal{C}'}$). Finally, if $G$ is a functor quasi-inverse of~$F$,
and if one chooses isomorphisms $\varphi\colon
FG\isomto\id_{\cal{C}'}$ and $\psi\colon GF\isomto\id_{\cal{C}}$, then
the two compatibility conditions on~$\varphi$, $\psi$
stated in \cite[I.1.2]{VI.1} are in fact equivalent to one
another, and for every chosen isomorphism $\varphi$, there exists
a unique isomorphism $\psi$ such that the said conditions are
satisfied.

\section{Categories over another}
\label{VI.2}

Let $\cal{E}$ be a category in~$\Univ$, it is thus an object
of~$\Cat$, and one can consider the category
$\Cat_{/\cal{E}}$
of ``objects of~$\Cat$ over~$\cal{E}$''. An object of this
category is thus a functor
$$
p\colon \cal{F}\to\cal{E}
$$
One also says that the category~$\cal{F}$, equipped with such a functor,
is a \emph{category over}~$\cal{E}$, or an
$\cal{E}$\emph{-category}. One will thus call an $\cal{E}$-functor
from a category $\cal{F}$ over~$\cal{E}$ to a category
$\cal{G}$ over~$\cal{E}$, a functor
$$
f\colon\cal{F}\to\cal{G}
$$
such that
$$
qf=p
$$
where $p$ and $q$ are the projection functors for $\cal{F}$
\resp $\cal{G}$. The set of $\cal{E}$-functors $f$ from~$\cal{F}$
to $\cal{G}$ is thus in one-to-one correspondence with the set
of arrows with source $\cal{F}$ and target $\cal{G}$
in~$\Cat_{/\cal{E}}$, without
% original p. 149
there being an identity there, however (since the datum of an
$f$ as above does not determine $\cal{F}$ and $\cal{G}$ as
categories over~$\cal{E}$); but of course, as in any
other category~$\cal{C}_{/S}$, one will commonly make the abuse of
language consisting in identifying the $\cal{E}$-functors (in the
sense made explicit above) with arrows in a
category~$\Cat_{/\cal{E}}$.

One will denote by
$$
\Hom_{\cal{E}}(\cal{F},\cal{G})
$$
the set of $\cal{E}$-functors from~$\cal{F}$ to~$\cal{G}$. Of
course, a composite of~$\cal{E}$-functors is an
$\cal{E}$-functor (the composition in question corresponding by
definition to the composition of arrows
in~$\Cat_{/\cal{E}}$).

Consider now two $\cal{E}$-functors
$$
f,g\colon \cal{F}\to\cal{G}
$$
and a homomorphism of functors:
$$
u\colon f\to g
$$
One says that $u$ is an $\cal{E}$-\emph{homomorphism} or a
``\emph{homomorphism of~$\cal{E}$-functors}'', if for every
$\xi\in\Ob(\cal{F})$, one has
$$
q(u(\xi))=\id_{p(\xi)} \quoi\text{,}
$$
in words: setting
$S=p(\xi)=qf(\xi)=qg(\xi)\in\Ob(\cal{E})$,
the morphism
$$
u(\xi)\colon f(\xi)\to g(\xi)
$$
in $\cal{G}$ is an $\id_S$-morphism. (In
general, for every morphism $\alpha\colon T \to S$
in~$\cal{E}$, and every category $\cal{G}$ over
$\cal{E}$, a morphism $v$ in $\cal{G}$ is called an
$\alpha$-\emph{morphism} if $q(v)=\alpha$, $q$ denoting the
projection functor $\cal{G}\to \cal{E}$). If one has a third
$\cal{E}$-functor $h\colon \cal{F}\to \cal{G}$ and an
$\cal{E}$-homomorphism $v\colon g\to h$, then $vu$ is likewise
an $\cal{E}$-homomorphism. Thus,
% original p. 150
\emph{the $\cal{E}$-functors from~$\cal{F}$ to~$\cal{G}$, and
the $\cal{E}$-homomorphisms of such, form a subcategory
of the category $\SheafHom(\cal{F},\cal{G})$ of all functors
from~$\cal{F}$ to~$\cal{G}$, which one will call the category of
$\cal{E}$-functors from~$\cal{F}$ to $\cal{G}$ and which one will denote}
$$
\SheafHom_{\cal{E}/-}(\cal{F},\cal{G})
$$
It is also the kernel subcategory of the pair of functors
$$
\xymatrix@C=.5cm{R,S\colon\SheafHom(\cal{F},\cal{G})
\ar@<2pt>[r]\ar@<-2pt>[r]&\SheafHom(\cal{F},\cal{E})}\quoi\text{,}
$$
where $R$ is the constant functor defined by the object $p$ of
$\SheafHom(\cal{F},\cal{E})$, and where $S$ is the functor $f\mto
q\circ f$ defined by $q\colon\cal{G}\to\cal{E}$.

To finish these generalities, it remains to define the
natural pairings among the categories
$\SheafHom_{\cal{E}/-}(\cal{F},\cal{G})$ by composition of
$\cal{E}$-functors. In other words, one wants to define a
``composition functor'':
\begin{equation*}
\tag{i}
\label{eq:VI.2.i} {\SheafHom_{\cal{E}/-}(\cal{F},\cal{G})\times
\SheafHom_{\cal{E}/-}(\cal{G},\cal{H})\to
\SheafHom_{\cal{E}/-}(\cal{F},\cal{H})}
\end{equation*}
when $\cal{F}$, $\cal{G}$, $\cal{H}$ are three categories
over~$\cal{E}$, in such a way that this functor induce, for the
objects, the composition map $(f,g)\mto gf$ of
$\cal{E}$-functors $f\colon\cal{F}\to\cal{G}$ and $g\colon \cal{G}\to
\cal{H}$. For this, let us recall that one defines a canonical
functor:
\begin{equation*}
\tag{ii}
\label{eq:VI.2.ii}
{\SheafHom(\cal{F},\cal{G})\times\SheafHom(\cal{G},\cal{H})\to
\SheafHom(\cal{F},\cal{H})}
\end{equation*}
which, for the objects, is nothing other than the composition map
$(f,g)\mto gf$ of functors, and which transforms an arrow
$(u,v)$,~where
$$
u\colon f\to f'\quoi,\quad v\colon g\to g'
$$
are arrows in $\SheafHom(\cal{F},\cal{G})$ \resp in
$\SheafHom(\cal{G},\cal{H})$, into the arrow
$$
v*u\colon gf\to g'f'
$$
defined by the relation:
% original p. 151
$$
v*u(\xi)=v(f'(\xi))\cdot g(u(\xi))=g'(u(\xi))\cdot v(f(\xi))
$$
It is well known that one thus indeed obtains a homomorphism from~$gf$
to~$g'f'$, and that (for $f,g$ and $u,v$ variable) one thus obtains
a functor~\eqref{eq:VI.2.ii}, \ie that one has
\begin{equation*}
\tag{I}
\label{eq:VI.2.I} {\id_g*\id_f=\id_{gf}\quoi,}
\end{equation*}
\begin{equation*}
\tag{II}
\label{eq:VI.2.II} {(v'*u')\circ(v*u)=(v'\circ v)*(u'\circ u)}
\end{equation*}
Let us also recall that one has an associativity formula for the
canonical pairings~\eqref{eq:VI.2.ii}, which is expressed on the one hand
by the associativity $(hg)f=h(gf)$ of composition of functors,
and on the other hand by the formula
\begin{equation*}
\tag{III}
\label{eq:VI.2.III} {(w*v)*u=w*(v*u)}
\end{equation*}
for the composition products of homomorphisms of functors (where
$u\colon f\to f'$ and $v\colon g\to g'$ are as above, and where one
assumes given moreover a homomorphism $w\colon h\to h'$ of
functors $h,h'\colon \cal{H}\to\cal{K}$). I now say that
\emph{when~$\cal{F}$, $\cal{G}$ are $\cal{E}$-categories,
the canonical composition functor}~\eqref{eq:VI.2.ii} \emph{induces a
functor}~\eqref{eq:VI.2.i}. As one already knows that the
composite of two $\cal{E}$-functors is an $\cal{E}$-functor,
this amounts to saying that \emph{when $u\colon f\to f'$ and
$v\colon g\to g'$ are homomorphisms of~$\cal{E}$-functors,
then $v*u\colon gf\to g'f'$ is likewise a homomorphism of
$\cal{E}$-functors.} This follows indeed trivially from the
definitions. As the pairings~\eqref{eq:VI.2.i} are
induced by the pairings~\eqref{eq:VI.2.ii}, they satisfy the
same associativity property, expressed also
in the formulae $(hg)f=h(gf)$ and $(w*v)*u=w*(v*u)$ for
$\cal{E}$-functors and $\cal{E}$-homomorphisms of
$\cal{E}$-functors.

To complete the formulary~\eqref{eq:VI.2.I},
\eqref{eq:VI.2.II}, \eqref{eq:VI.2.III}, let us also recall the formulae:
\begin{equation*}
\tag{IV}
\label{eq:VI.2.IV} {v*\id_{\cal{F}}=v\quad\text{and}\quad\id_{\cal{G}}*u=u}\quoi,
\end{equation*}
where
% original p. 152
for simplicity one writes $v*f$ or $u*g$ instead of~$v*u$, when
$u$ \resp $v$ is the identity automorphism of~$f$ \resp $g$.

From the definition of the pairings~\eqref{eq:VI.2.i} it follows
that $\SheafHom_{\cal{E}/-}(\cal{F},\cal{G})$ \emph{is a functor in
$\cal{F}$,~$\cal{G}$,
of the product category
${\Cat_{/\cal{E}}}^\circ\times\Cat_{/\cal{E}}$ to the
category~$\Cat$}. If indeed $g\colon \cal{G}\to\cal{G}_1$ is
an $\cal{E}$-functor, \ie an object of
$\SheafHom_{\cal{E}/-}(\cal{G},\cal{G}_1)$, then setting
in~\eqref{eq:VI.2.i} $\cal{H}=\cal{G}_1$, there corresponds to it a
functor
$$
g_*\colon\SheafHom_{\cal{E}/-}(\cal{F},\cal{G})\to
\SheafHom_{\cal{E}/-}(\cal{F},\cal{G}_1)
$$
One defines analogously, for an $\cal{E}$-functor
$f\colon \cal{F}_1\to \cal{F}$, a functor
$$
f^*:\SheafHom_{\cal{E}/-}(\cal{F},\cal{G})\to
\SheafHom_{\cal{E}/-}(\cal{F}_1,\cal{G})
$$
For brevity, one also denotes these functors by the symbols
$f\mto g\circ f$ \resp $g\mto g\circ f$ (which designate
only, in fact, the corresponding maps on the sets
of objects). It follows from the associativity property
indicated above that in this way, one indeed obtains as
announced a functor
$\Cat_{/\cal{E}}^\circ\times\Cat_{/\cal{E}}\to\Cat$.

\section{Change of base in categories over~$\cal{E}$}
\label{VI.3}

As in $\Cat$ the projective limits (relative to
categories $\cal{I}$ elements of~$\Univ$) exist, the same
is true in $\Cat_{/\cal{E}}$, in particular cartesian
products exist there, which are interpreted as fibered
products in~$\Cat$. In accordance with the general
notations, if $\cal{F}$ and $\cal{G}$ are categories
over~$\cal{E}$, one denotes by
$$
\cal{F}\times_{\cal{E}}\cal{G}
$$
their product in~$\Cat_{/\cal{E}}$, \ie their fibered product
over~$\cal{E}$ in~$\Cat$,
% original p. 153
as a category over~$\cal{E}$. Thus,
$\cal{F}\times_{\cal{E}} \cal{G}$ is equipped with two $\cal{E}$-functors
$\pr_1$ and $\pr_2$, which define, for every category
$\cal{H}$ over~$\cal{E}$, a bijection
$$
\Hom_{\cal{E}}(\cal{H},\cal{F}\times_{\cal{E}}\cal{G})\isomto
\Hom_{\cal{E}}(\cal{H},\cal{F})\times\Hom_{\cal{E}}(\cal{H},\cal{G}).
$$
This bijection arises moreover from an isomorphism of
categories
$$
\SheafHom_{\cal{E}/-}(\cal{H},\cal{F}\times_{\cal{E}}
\cal{G})\isomto\SheafHom_{\cal{E}/-}(\cal{H},\cal{F})\times
\SheafHom_{\cal{E}/-}(\cal{H},\cal{G})
$$
by taking the sets of objects of the two sides, where the functor
written is the one which has as components the functors $h\mto
\pr_1\circ h$ and $h\mto \pr_2\circ h$ from the first side to the two
factors of the second. The reader is left to verify that one
thus indeed obtains an isomorphism (the analogous fact being true,
more generally, whenever one has a projective limit
of categories --- and not only in the case of a fibered
product).

Let us moreover recall that one has (as was said in~no.~\ref{VI.1}):
\begin{eqnarray*}
\Ob(\cal{F}\times_\cal{E}\cal{G})&=&
\Ob(\cal{F})\times_{\Ob(\cal{E})}\Ob(\cal{G})\\
\Fl(\cal{F}\times_\cal{E}\cal{G})&=&
\Fl(\cal{F})\times_{\Fl(\cal{E})}\Fl(\cal{G})\quoi,
\end{eqnarray*}
the composition of arrows being performed moreover
componentwise.

In the sequel, we consider a functor
$$
\lambda\colon\cal{E}'\to\cal{E}
$$
and for every category $\cal{F}$ over~$\cal{E}$, one
considers $\cal{F}\times_\cal{E}\cal{E}'$ as a category
over~$\cal{E}'$ by means of $\pr_2$; in other words,
we interpret the operation ``fibered product'' as an
operation \emph{``change of base''}, the functor
$\lambda\colon\cal{E}'\to\cal{E}$ taking the name of \emph{``change of
base functor''.}
In accordance with the well-known general facts, one thus obtains
a functor, called the
% original p. 154
\emph{change of base functor} for $\lambda$:
$$
\lambda^*\colon\Cat_{/\cal{E}}\to\Cat_{/\cal{E}'}\quoi,
$$
(adjoint of the functor ``restriction of the base'' which, to every
category $\cal{F}'$ over~$\cal{E}'$, with
projection functor~$p'$, associates~$\cal{F}'$, regarded as a
category over~$\cal{E}$ by the functor~$p=\lambda
p'$). As is well known in the general case of a
change of base functor in a category, the
change of base functor ``commutes with projective limits'', and in
particular ``transforms'' fibered products over~$\cal{E}$ into
fibered products over~$\cal{E}'$.

Let $\cal{F}$ and $\cal{G}$ be two categories over
$\cal{E}$, one will define a \emph{canonical isomorphism}:
\begin{multline*}
\tag{i}
\label{eq:VI.3.i}
\SheafHom_{\cal{E}'/-}(\cal{F}',\cal{G}')\isomto
\SheafHom_{\cal{E}/-}(\cal{F}\times_{\cal{E}}\cal{E}',\cal{G})\\
\text{where }\cal{F}'=\cal{F}\times_{\cal{E}}\cal{E}'\text{, }
\cal{G}'=\cal{G}\times_{\cal{E}}\cal{E}'.
\end{multline*}
For this, consider the functor
$$
\pr_1\colon\cal{G}'=\cal{G}\times_{\cal{E}}\cal{E}'\to \cal{G}\quoi,
$$
and define~\eqref{eq:VI.3.i} by
$$
F\mto \pr_1\circ F\quoi,
$$
which a priori designates a functor
\begin{equation*}
\tag{ii}
\label{eq:VI.3.ii} {\SheafHom(\cal{F}',\cal{G}')\to\SheafHom(\cal{F}',\cal{G})}
\end{equation*}
It is therefore only necessary to verify that the latter induces a functor
for the subcategories~\eqref{eq:VI.3.i}, and that the latter is
an isomorphism. That~\eqref{eq:VI.3.ii} induces a bijection
$$
\Hom_{\cal{E}'/-}(\cal{F}',\cal{G}')\isomto
\Hom_{\cal{E}/-}(\cal{F}\times_{\cal{E}}\cal{E}',\cal{G})
$$
is the characteristic property of the change of
base functor. It remains therefore to
% original p. 155
prove that if $F$, $G$ are $\cal{E}'$-functors
$\cal{F}'\to\cal{G}'$, then \emph{the map}
$$
u\mto \pr_1\circ u
$$
\emph{induces a bijection}
$$
\Hom_{\cal{E}'}(F,G)\isomto\Hom_{\cal{E}}(\pr_1\circ F,\pr_1\circ
G)\quoi.
$$
The verification of this fact is immediate, and left to the
reader.

It follows from this isomorphism~\eqref{eq:VI.3.i}, and from the end of
the preceding no., that
$$
\SheafHom_{\cal{E}'/-}(\cal{F}\times_{\cal{E}}\cal{E}',\cal{G}
\times_{\cal{E}}\cal{E}')
$$
\emph{can be regarded as a functor in}
$\cal{E}',\cal{F},\cal{G}$, \emph{from the category}
$\Cat_{/\cal{E}}^\circ\times\Cat_{/\cal{E}}^\circ\allowbreak\times\Cat_{/\cal{E}}$
\emph{to the category}~$\Cat$, isomorphic to the functor defined
by the expression
$\SheafHom_{\cal{E}/-}(\cal{F}\times_{\cal{E}}\cal{E}',\cal{G})$. In
particular, for $\cal{F}$,~$\cal{G}$ fixed, one obtains a
functor in $\cal{E}'$, $\cal{E}'\to
\SheafHom_{\cal{E}''/-}(\cal{F}',\cal{G}')=
\SheafHom_{\cal{E}'/-}(\cal{F}\times_{\cal{E}}
\cal{E}',\cal{G}\times_{\cal{E}}\cal{E}')$, and in particular the
projection $\cal{E}$-functor $\lambda\colon \cal{E}'\to\cal{E}$
defines a morphism \ie a functor
$$
\lambda^*_{\cal{F},\cal{G}}\colon
\SheafHom_{\cal{E}/-}(\cal{F},\cal{G})\to
\SheafHom_{\cal{E}'/-}(\cal{F}',\cal{G}')
$$
which we shall make explicit. For the sets of objects of the two
sides, it is the map
$$
f\mto f'=f\times_{\cal{E}}\cal{E}'
$$
which expresses the functorial dependence of
$\cal{F}\times_{\cal{E}}\cal{E}'$ on the object $\cal{F}$
over~$\cal{E}$. On the other hand, consider two $\cal{E}$-functors
$$
f,g\colon\cal{F}\to \cal{G}
$$
and a homomorphism of~$\cal{E}$-functors
$$
u\colon f\to g\quoi,
$$
we
% original p. 156
shall make explicit the corresponding homomorphism of~$\cal{E}'$-functors:
$$
u'\colon f'\to g'\quoi.
$$
For every
$$
\xi'=(\xi,S')\in\Ob(\cal{F}')
$$
with
$$
\xi\in\Ob(\cal{F}),\quad S'\in\Ob(\cal{E}'),\quad p(\xi)=\lambda(S')=S
$$
the morphism
$$
u'(\xi')\colon f'(\xi')=(f(\xi),S')\to g'(\xi')=(g(\xi),S')\quad
\text{in } \cal{G}'
$$
is defined by the formula
$$
u'(\xi')=(u(\xi),\id_{S'})
$$
(which is indeed an $S'$-morphism in~$\cal{G}'$, because
$q(u(\xi))=\lambda(\id_{S'})=\id_S$).

Consider now an arbitrary $\cal{E}$-functor
$$
\lambda'\colon\cal{E}''\to\cal{E}'
$$
and the corresponding functor
$$
\SheafHom_{\cal{E}'/-}(\cal{F}\times_{\cal{E}}\cal{E}',
\cal{G}\times_{\cal{E}}\cal{E}')\to
\SheafHom_{\cal{E}''/-}(\cal{F}\times_{\cal{E}}\cal{E}'',
\cal{G}\times_{\cal{E}}\cal{E}'')\quoi,
$$
I say that this functor is nothing other than the functor one obtains by
the preceding procedure, starting from~$\cal{F}'$ and
$\cal{G}'$ over~$\cal{E}'$ and regarding $\cal{E}''$ as a
category over~$\cal{E}'$, taking into account the isomorphisms of
\emph{``transitivity of change of base''}:
$$
\cal{F}'\times_{\cal{E}'}\cal{E}''\isomto
\cal{F}''=\cal{F}\times_{\cal{E}}\cal{E}''\quad
\text{and}\quad\cal{G}'\times_{\cal{E}'}\cal{E}''\isomto\cal{G}''
=\cal G\times_{\cal E}\cal E''\quoi,
$$
which imply a canonical isomorphism
$$
\SheafHom_{\cal{E}''/-}(\cal{F}'\times_{\cal{E}'}\cal{E}'',
\cal{G}'\times_{\cal{E}'}\cal{E}'')\isomto
\SheafHom_{\cal{E}''/-}(\cal{F}\times_{\cal{E}}\cal{E}'',
\cal{G}\times_{\cal{E}}\cal{E}'')
$$
The
% original p. 157
verification of this compatibility is immediate, and
left to the reader.

The functors just defined are compatible with the
pairings defined in the preceding no., in a precise way: if $\cal{F}$, $\cal{G}$,~$\cal{H}$ are
categories over~$\cal{E}$ and if one sets
$$
\cal{F}'=\cal{F}\times_{\cal{E}}\cal{E}'\quoi,\quad \cal{G}'=
\cal{G}\times_{\cal{E}}\cal{E}'\quoi,\quad \cal{H}'=
\cal{H}\times_{\cal{E}}\cal{E}'\quoi,
$$
one has commutativity in the following diagram of functors:
$$
\xymatrix{ \SheafHom_{\cal{E}/-}(\cal{F},\cal{G})\times
\SheafHom_{\cal{E}/-}(\cal{G},\cal{H}) \ar[r]
\ar[d]_{\lambda^*_{\cal{F},\cal{G}}\times \lambda^*_{\cal{G},\cal{H}}}
& \SheafHom_{\cal{E}/-}(\cal{F},\cal{H})
\ar[d]^{\lambda^*_{\cal{F},\cal{H}}} \\
\SheafHom_{\cal{E}'/-}(\cal{F}',\cal{G}')\times
\SheafHom_{\cal{E}'/-}(\cal{G}',\cal{H}')\ar[r] &
\SheafHom_{\cal{E}'/-}(\cal{F}',\cal{H}')}
$$
where the horizontal arrows are the composition functors
defined in the preceding no. This commutativity
is expressed by the formulae
$$
(gf)'=g'f'
$$
for $f\in\Hom_{\cal{E}}(\cal{F},\cal{G})$,
$g\in\Hom_{\cal{E}}(\cal{G},\cal{H})$, (a formula which simply expresses
the functoriality of change of base), and the formula
$$
(v*u)'=v'*u'
$$
when $u\colon f\to f_1$ is an arrow of
$\SheafHom_{\cal{E}/-}(\cal{F},\cal{G})$ and $v\colon g\to g_1$ an
arrow of~$\SheafHom_{\cal{E}/-}(\cal{G},\cal{H})$. The
verification of this formula follows easily from the
definitions.

In the sequel, we shall be interested mainly in
$\SheafHom_{\cal{E}}(\cal{F},\cal{G})$ (and certain
remarkable subcategories of the latter) when
$\cal{F}=\cal{E}$, and we introduce for this reason a
special notation:
$$
\mathbf{\Gamma}(\cal{G}/\cal{E})=
\SheafHom_{\cal{E}}(\cal{E},\cal{G})\quoi,\quad
\Gamma(\cal{G}/\cal{E})=\Ob(\mathbf{\Gamma}(\cal{G}/\cal{E}))=
\Hom_{\cal{E}}(\cal{E},\cal{G})\quoi.
$$

\begin{remarks}
When
% original p. 158
$\cal{E}$ is a punctual category,
\ie $\Ob(\cal{E})$ and $\Fl(\cal{E})$ reduced to a single
element, which also means that $\cal{E}$ is a
final object of the category~$\Cat$, then the datum of a
category over~$\cal{E}$ is equivalent to the datum
of just a category, (because there will be a unique functor from
$\cal{F}$ to~$\cal{E}$). More precisely,
$\Cat_{/\cal{E}}$ is then isomorphic to~$\Cat$. Moreover, the
categories $\SheafHom_{\cal{E}/-}(\cal{F},\cal{G})$ are then nothing
other than the $\SheafHom(\cal{F},\cal{G})$. Let us then recall that the
fundamental formula
$$
\Hom(\cal{H},\SheafHom(\cal{F},\cal{G}))\isomto
\Hom(\cal{F}\times\cal{H},\cal{G})
$$
(a functorial isomorphism in the three arguments which appear there),
allows one to interpret $\SheafHom(\cal{F},\cal{G})$ axiomatically,
in terms internal to the category~$\Cat$, so that the
known formulary of the $\SheafHom$ categories appears as a
special case of a formulary valid in categories such
as~$\Cat$, where ``$\SheafHom$-objects'' (defined by the
preceding formula) exist. There is an analogous
interpretation of~$\SheafHom_{\cal{E}/-}(\cal{F},\cal{G})$ when one assumes
once more that $\cal{E}$ is arbitrary, by the formula
$$
\Hom(\cal{H},\SheafHom_{\cal{E}/-}(\cal{F},\cal{G}))\isomto
\Hom_{\cal{E}}(\cal{F}\times\cal{H},\cal{G})
$$
(a functorial isomorphism in the three arguments). In this way,
the formal properties set out in~no.~\ref{VI.2},~\ref{VI.3} are
special cases of more general results, valid
in categories where the objects
$\SheafHom_{\cal{E}/-}(\cal{F},\cal{G})$ (when $\cal{F}$,~$\cal{G}$
are two objects of the category over a third~$\cal{E}$) exist.
\end{remarks}

% ---- en-chunk2.tex ----
